\documentclass[11pt]{article}
\usepackage[margin=1.1in]{geometry}
\usepackage{amsmath,amssymb,amsthm,mathtools,booktabs,array,tabularx,longtable}
\usepackage[dvipsnames]{xcolor}
\usepackage[colorlinks=true,linkcolor=MidnightBlue,citecolor=MidnightBlue,urlcolor=MidnightBlue]{hyperref}
\usepackage{fancyhdr}
\pagestyle{fancy}\fancyhf{}\renewcommand{\headrulewidth}{0pt}
\fancyfoot[R]{\thepage}
\fancyfoot[L]{\scriptsize\color{gray}Generated by the Transport Lab from approved parts of its store; not a deposited paper. Papers of record: \href{https://zenodo.org/communities/transport-papers}{Zenodo, Transport Papers}}
\fancypagestyle{plain}{\fancyhf{}\fancyfoot[R]{\thepage}\fancyfoot[L]{\scriptsize\color{gray}Generated by the Transport Lab from approved parts of its store; not a deposited paper. Papers of record: \href{https://zenodo.org/communities/transport-papers}{Zenodo, Transport Papers}}\renewcommand{\headrulewidth}{0pt}}
\newtheorem{theorem}{Theorem}[section]
\newtheorem{proposition}[theorem]{Proposition}
\newtheorem{lemma}[theorem]{Lemma}
\newtheorem{corollary}[theorem]{Corollary}
\newtheorem{conjecture}[theorem]{Conjecture}
\newtheorem{refutation}[theorem]{Refutation}
\newtheorem{prediction}[theorem]{Prediction}
\newtheorem{principle}[theorem]{Principle}
\newtheorem{observation}[theorem]{Observation}
\newtheorem{axiom}[theorem]{Axiom}
\theoremstyle{definition}
\newtheorem{definition}[theorem]{Definition}
\newtheorem{fact}[theorem]{Fact}
\newtheorem{assembled}[theorem]{Assembly}
\newtheorem{postulate}[theorem]{Postulate}
\newtheorem{premise}[theorem]{Premise}
\newtheorem{hypothesis}[theorem]{Hypothesis}
\newtheorem{remark}[theorem]{Remark}
\newtheorem{example}[theorem]{Example}
\theoremstyle{plain}
\setlength{\marginparwidth}{0.8in}
\newcommand{\changedmark}[2]{\marginpar{\raggedright\scriptsize\color{#1}#2}}
\providecommand{\Aa}{}\renewcommand{\Aa}{\mathcal{A}}
\providecommand{\Ad}{}\renewcommand{\Ad}{\widetilde{\mathrm{Ad}}}
\providecommand{\Aev}{}\renewcommand{\Aev}{\mathcal{A}^{+}}
\providecommand{\Aod}{}\renewcommand{\Aod}{\mathcal{A}^{-}}
\providecommand{\B}{}\renewcommand{\B}{\mathrm{B}}
\providecommand{\C}{}\renewcommand{\C}{\mathbb{C}}
\providecommand{\CAR}{}\renewcommand{\CAR}{\mathrm{CAR}}
\providecommand{\CS}{}\renewcommand{\CS}{{\mathbb{S}^{+}}}
\providecommand{\CSm}{}\renewcommand{\CSm}{{\mathbb{S}^{-}}}
\providecommand{\Cl}{}\renewcommand{\Cl}{\mathrm{Cl}}
\providecommand{\Con}{}\renewcommand{\Con}{\mathcal{C}}
\providecommand{\Dl}{}\renewcommand{\Dl}{\Delta}
\providecommand{\EN}{}\renewcommand{\EN}{E_{N}}
\providecommand{\End}{}\renewcommand{\End}{\mathrm{End}}
\providecommand{\F}{}\renewcommand{\F}{\mathbb{F}}
\providecommand{\G}{}\renewcommand{\G}{\Gamma}
\providecommand{\HH}{}\renewcommand{\HH}{\mathbb{H}}
\providecommand{\IS}{}\renewcommand{\IS}{I\sigma_3}
\providecommand{\J}{}\renewcommand{\J}{\mathbf{J}}
\providecommand{\LH}{}\renewcommand{\LH}{L_{H}}
\providecommand{\Mat}{}\renewcommand{\Mat}{\mathrm{M}}
\providecommand{\Mo}{}\renewcommand{\Mo}{M^{\circ}}
\providecommand{\Mp}{}\renewcommand{\Mp}{M^{+}}
\providecommand{\Msun}{}\renewcommand{\Msun}{M_{\odot}}
\providecommand{\Mt}{}\renewcommand{\Mt}{\widetilde{M}}
\providecommand{\Ncone}{}\renewcommand{\Ncone}{\mathcal{N}}
\providecommand{\Neff}{}\renewcommand{\Neff}{N_{\mathrm{eff}}}
\providecommand{\PCS}{}\renewcommand{\PCS}{\mathbb{P}{\mathbb{S}^{+}}}
\providecommand{\PT}{}\renewcommand{\PT}{\mathbb{PT}}
\providecommand{\Pin}{}\renewcommand{\Pin}{\mathrm{Pin}}
\providecommand{\Q}{}\renewcommand{\Q}{\mathbb{Q}}
\providecommand{\R}{}\renewcommand{\R}{\mathbb{R}}
\providecommand{\Ric}{}\renewcommand{\Ric}{\mathrm{Ric}}
\providecommand{\SBH}{}\renewcommand{\SBH}{S_{\mathrm{BH}}}
\providecommand{\Sc}{}\renewcommand{\Sc}{\mathbb{S}}
\providecommand{\Sch}{}\renewcommand{\Sch}{\mathrm{P}}
\providecommand{\Sig}{}\renewcommand{\Sig}{\Sigma}
\providecommand{\Som}{}\renewcommand{\Som}{S_\omega}
\providecommand{\Sp}{}\renewcommand{\Sp}{\mathbb{S}}
\providecommand{\Sph}{}\renewcommand{\Sph}{\mathrm{S}}
\providecommand{\Spi}{}\renewcommand{\Spi}{S_\pi}
\providecommand{\Spin}{}\renewcommand{\Spin}{\mathrm{Spin}}
\providecommand{\Sth}{}\renewcommand{\Sth}{S_{\mathrm{th}}}
\providecommand{\T}{}\renewcommand{\T}{\mathbb{T}}
\providecommand{\TH}{}\renewcommand{\TH}{T_{\mathrm{H}}}
\providecommand{\Tr}{}\renewcommand{\Tr}{\operatorname{Tr}}
\providecommand{\UHF}{}\renewcommand{\UHF}{\mathrm{UHF}}
\providecommand{\Z}{}\renewcommand{\Z}{\mathbb{Z}}
\providecommand{\Zc}{}\renewcommand{\Zc}{\mathrm{Z}}
\providecommand{\cmg}{}\renewcommand{\cmg}{\mathrm{cm^{2}\,g^{-1}}}
\providecommand{\dS}{}\renewcommand{\dS}{\mathrm{dS}}
\providecommand{\dsc}{}\renewcommand{\dsc}{\mathrm{dsc}}
\providecommand{\dtPG}{}\renewcommand{\dtPG}{dt}
\providecommand{\dtS}{}\renewcommand{\dtS}{dt_{\mathrm{S}}}
\providecommand{\ellv}{}\renewcommand{\ellv}{\ell}
\providecommand{\er}{}\renewcommand{\er}{e_r}
\providecommand{\et}{}\renewcommand{\et}{e_t}
\providecommand{\ev}{}\renewcommand{\ev}{\mathrm{ev}}
\providecommand{\gl}{}\renewcommand{\gl}{\mathfrak{gl}}
\providecommand{\hatV}{}\renewcommand{\hatV}{\widehat{V}}
\providecommand{\hb}{}\renewcommand{\hb}{\bar{h}}
\providecommand{\hf}{}\renewcommand{\hf}{\hat{f}}
\providecommand{\id}{}\renewcommand{\id}{\mathrm{id}}
\providecommand{\kpc}{}\renewcommand{\kpc}{\mathrm{kpc}}
\providecommand{\lP}{}\renewcommand{\lP}{\ell_{P}}
\providecommand{\mD}{}\renewcommand{\mD}{\mathbf{m}_{D}}
\providecommand{\mb}{}\renewcommand{\mb}{\bar m}
\providecommand{\mnu}{}\renewcommand{\mnu}{\Sigma m_{\nu}}
\providecommand{\nInf}{}\renewcommand{\nInf}{n_\infty}
\providecommand{\nO}{}\renewcommand{\nO}{n_0}
\providecommand{\ninf}{}\renewcommand{\ninf}{n_{\infty}}
\providecommand{\nz}{}\renewcommand{\nz}{n_{0}}
\providecommand{\proj}{}\renewcommand{\proj}{\operatorname{proj}}
\providecommand{\pv}{}\renewcommand{\pv}{\pi}
\providecommand{\rad}{}\renewcommand{\rad}{\mathrm{rad}}
\providecommand{\sgn}{}\renewcommand{\sgn}{\operatorname{sgn}}
\providecommand{\shA}{}\renewcommand{\shA}{\mathcal{A}}
\providecommand{\shC}{}\renewcommand{\shC}{\mathcal{C}}
\providecommand{\shE}{}\renewcommand{\shE}{\mathcal{E}}
\providecommand{\shJ}{}\renewcommand{\shJ}{\mathcal{J}}
\providecommand{\shK}{}\renewcommand{\shK}{\mathcal{K}}
\providecommand{\shL}{}\renewcommand{\shL}{L}
\providecommand{\shW}{}\renewcommand{\shW}{\mathcal{W}}
\providecommand{\sll}{}\renewcommand{\sll}{\mathfrak{sl}}
\providecommand{\so}{}\renewcommand{\so}{\mathfrak{so}}
\providecommand{\su}{}\renewcommand{\su}{\mathfrak{su}}
\providecommand{\supp}{}\renewcommand{\supp}{\operatorname{supp}}
\providecommand{\tr}{}\renewcommand{\tr}{\operatorname{tr}}
\providecommand{\uu}{}\renewcommand{\uu}{\mathfrak{u}}
\providecommand{\br}{}\renewcommand{\br}[1]{\langle #1\rangle}
\providecommand{\bra}{}\renewcommand{\bra}[1]{\langle #1|}
\providecommand{\file}{}\renewcommand{\file}[1]{\texttt{#1}}
\providecommand{\hp}{}\renewcommand{\hp}[2]{( #1 \,|\, #2 )}
\providecommand{\ip}{}\renewcommand{\ip}[2]{\langle #1, #2\rangle}
\providecommand{\ket}{}\renewcommand{\ket}[1]{|#1\rangle}
\providecommand{\rel}{}\renewcommand{\rel}[2]{D\bigl(#1 \,\|\, #2\bigr)}
\providecommand{\rev}{}\renewcommand{\rev}[1]{\widetilde{#1}}
\providecommand{\slashed}{}\renewcommand{\slashed}[1]{{#1}\mkern-10mu/}
\title{\bfseries Families of quarks and leptons}
\author{Leon Butler\\ \normalsize\href{mailto:lbutler.math@gmail.com}{\texttt{lbutler.math@gmail.com}}\\ \normalsize ORCID: \href{https://orcid.org/0009-0006-4449-5617}{0009-0006-4449-5617}}
\date{Generated by the Transport Lab, 2 November 2026}
\begin{document}
\maketitle
\begin{abstract}
\normalsize
Matter comes in three copies, or generations, of each kind of particle, with no accepted explanation for why there are three. This trail follows the conjecture that three counts the ways of choosing two building blocks out of three, and it ends with a plain answer: the series does not derive three generations.

Anomalies cancel within each generation, so consistency allows any number, and requiring the strong and weak forces to weaken at high energy allows at most five with one or two Higgs doublets. Three rests on an imported, observation-level bound read through a family-symmetry hypothesis.

The walk then asks how far the framework gets in telling the copies apart. The vacuum gives one heavy family and two light ones, but the simplest vacua leave the light pair exactly degenerate. The electron-to-muon ratio of about $207$ remains the open problem.
\end{abstract}
\newpage
\noindent \textbf{Authorship and disclosure. }The founding idea of the series is the author's: that division by zero in a Clifford algebra should be read not as a prohibition but as a transport into a higher-dimensional algebra, and that the transported division need not be accompanied by an inverse operation. The direction of investigation and the critical review of the results are likewise the author's. The mathematical development and the drafting of the text were carried out in extended collaboration with Claude Fable 5 (Anthropic), an AI system.

\noindent\fbox{\parbox{\dimexpr\linewidth-2\fboxsep-2\fboxrule}{\textbf{This is not a deposited paper.} It was generated by the Transport Lab on 2 November 2026 and assembled from approved parts of its store. It follows results across several papers of the series; each is stated with its proof and credited to the paper that first proved it. The official versions of those papers are on Zenodo, in the Transport Papers community, \url{https://zenodo.org/communities/transport-papers}.}}
\newpage

\section*{Introduction}
The walk has two parts. First it asks whether the framework can count the families, starting from two facts about quantum fields that bear on the count without fixing it. Then it follows the attempts to make identical families differ in mass, ending where the account stops.

Negative results stay on the trail, because each one narrows what the answer can be.

\section{What we assume}

\begin{postulate}[The slot-module hypothesis: the weaker constituent hypothesis behind (S), with the slot module generated from a unique vacuum and the chirality a CPT halving commuting with the Pati-Salam algebra]\label{c:standard-model.slot-module-hypothesis}
The framework's account of quarks and leptons needed an assumption, called (S), that two particular ingredients of each particle behave like fermions. This weaker assumption replaces it. It says how a particle's possible states are organised, that they are built from exactly one empty state by adding quanta of a single two-state kind, that handedness splits the states into particles and their mirror-image antiparticles in the way every theory of fermions requires, with that split respecting the intermediate Pati-Salam symmetry the framework already uses, and that, if the ingredients are fermions, they combine with the particle's other fermionic modes into one system. Those requirements, with the standard consistency condition that the theory be free of anomalies, force the ingredients to behave like fermions and fix which states are left-handed, so (S) becomes a consequence rather than an assumption. Respecting the Pati-Salam symmetry does one job the others cannot: it decides that the completely empty state counts as a particle rather than an antiparticle.
\end{postulate}

\begin{postulate}[(KF), the Krein-form premise for gauge fields: on the palatial storey the gauge fields, their ghosts and their Pauli-Villars regulator fields carry one indefinite (Krein) inner product, the Gupta-Bleuler form $[a(u), a(w)^{\dagger}] = -u\cdot w$ on the polarisations, and the physical space is positive through the quartet mechanism]\label{c:synthesis.krein-form-premise-indefinite-inner-product-for-gauge-fields}
To describe light and the other force fields in a way that treats space and time alike, physicists use bookkeeping in which some unphysical states have negative or zero length, and remove them again before anything is measured. The same device is used to tame the infinities of quantum corrections, with auxiliary particles that count negatively. The series' quantum spaces all have positive lengths, so it cannot use this bookkeeping as it stands. This premise adopts it: the force fields of the series carry the standard indefinite measure built from the spacetime metric of their polarisations, and only states of positive length remain once the unphysical parts cancel in groups of four. Several results use it, among them the series' calculations of the first and second quantum corrections of electrodynamics, its calculation of how the force strengths change with energy, and its construction of the ghost bookkeeping. One of the two uses, the auxiliary negative particles, turns out not to be needed for the screening of charge by the vacuum, which can be computed without them. It is taken as given, not derived; it would be refuted if the physical states it leaves were found to include ones of negative length.
\end{postulate}

\section{Two checks on the count}

The walk starts with two standard facts of quantum field theory, recomputed in the spacetime algebra, because each depends on the number of generations in a different way.

Matter comes in three generations, and nobody has an accepted explanation for why there are three. A first place to look for the number is consistency.

A quantum field theory with chiral forces must have its triangle anomalies cancel. In the spacetime algebra the anomaly takes a clean form: the axial divergence is the bilinear form of $\langle I_{4}F^{2}\rangle_{0} = -2E\cdot B$, the pseudoscalar part of the square of the field bivector. This is the standard result of Adler, Bell and Jackiw, recomputed here with imported Pauli-Villars regulators.

\begin{theorem}[The triangle anomaly in the spacetime algebra with imported Pauli-Villars regulators (standard: Adler; Bell and Jackiw): the vector Ward identities hold integrand by integrand, the axial divergence is the regulator-mass limit $(1/(2\pi^{2}))\langle I_{4}\gamma^{\nu}k_{1}\gamma^{\rho}k_{2}\rangle_{0}$ per Dirac fermion for any weights with $\sum_{i} c_{i} = 0$, the bilinear form of the pseudoscalar part of $F^{2}$, and for one stored generation the gauge anomaly tensor vanishes on the whole gauge algebra (standard per-generation cancellation)]\label{c:synthesis.triangle-anomaly-in-the-spacetime-algebra-is-the-pseudoscalar-part-of-the-field-square-and-cancels-per-generation}
Some symmetries of classical fields fail once quantum effects are included; the classic case is the axial anomaly of a charged particle in an electromagnetic field, found in 1969. This result recomputes that anomaly in the series' own algebra of spacetime, using a standard method of taming the infinities. The anomaly comes out as the part of the squared electromagnetic field that measures how electric and magnetic fields line up, with the textbook strength. For a theory to be consistent, such anomalies must cancel between all the particles that feel a force; the result checks that they cancel within one generation of the series' particles, so the cancellation does not depend on how many generations there are.
\end{theorem}
\noindent\emph{Rests on:} \ref{c:synthesis.one-loop-qed-in-the-spacetime-algebra}, \ref{c:synthesis.imported-one-loop-integrals-and-pauli-villars-regulators}, \ref{c:synthesis.anomaly-freedom-forces-slot-content}, \ref{c:synthesis.imported-adler-bell-jackiw-axial-anomaly}.

The anomaly tensor vanishes within one generation of the stored content, so it cancels generation by generation and says nothing about how many generations there are.

A second place to look is how the strengths of the forces change with energy. Spin traces in the spacetime algebra, turned into running weights by an imported standard heat-kernel formula, give $\tfrac13$, $\tfrac23$ and $-\tfrac{11}{3}$ for complex scalars, Weyl spinors and gauge fields with their ghosts. The computation works in a covariant gauge, where the gauge fields carry an indefinite (Krein) inner product, adopted as a premise.

With $n_{g}$ generations and $n_{h}$ Higgs doublets as inputs, three generations and one doublet give the standard coefficients $(\tfrac{41}{10}, -\tfrac{19}{6}, -7)$.

\begin{theorem}[The one-loop running coefficients from spin traces in the spacetime algebra, in a covariant gauge under the indefinite-form premise: the Lorentz generator of the field-strength bivector gives spin weights $0$, $-1$ and $-4$ for scalars, Weyl spinors (a minimal left ideal with the unit $I_{4}$) and vectors, the imported heat-kernel coefficient turns them into $\tfrac13$, $\tfrac23$ and $-\tfrac{11}{3}$ per unit index, and under (S), with $n_{g}$ generations and $n_{h}$ doublets as inputs, the coefficients are $(\tfrac{4n_{g}}{3} + \tfrac{n_{h}}{10}, -\tfrac{22}{3} + \tfrac{4n_{g}}{3} + \tfrac{n_{h}}{6}, -11 + \tfrac{4n_{g}}{3})$, asymptotically free in both non-abelian factors for $n_{g} = 3$ and every $n_{h} \leq 19$]\label{c:synthesis.one-loop-running-coefficients-from-spacetime-algebra-spin-traces}
How fast the strengths of the forces change with energy is fixed, at the first quantum correction, by a sum over the particles, each weighted by a number that depends on its spin. Those weights come out of the series' algebra of spacetime as traces of how a field's spin responds to a force field, once a standard heat-kernel formula and the standard covariant treatment of the force fields are taken from the literature. The number of generations and of Higgs doublets are inputs, since the store has not settled either. Applied to the particle content the series derives under its slot hypothesis, three generations with one doublet give the standard values and two doublets give slightly different ones; with three generations the strong and weak forces weaken at high energy for any number of doublets up to nineteen. The gluons and the weak bosons are not yet built as quantum states here.
\end{theorem}
\noindent\emph{Rests on:} \ref{c:field-theory.second-quantisation}, \ref{c:field-theory.fock-space-maxwell-modes}, \ref{c:standard-model.gauge-algebra-minimal-massive-composite}, \ref{c:standard-model.hypercharge-global-form-pair-counting-consistency}, \ref{c:standard-model.imported-one-loop-running-trace-ratio}, \ref{c:synthesis.imported-heat-kernel-coefficients-and-background-field-gauge}, \ref{c:synthesis.krein-form-premise-indefinite-inner-product-for-gauge-fields}.

\section{The faces as families}

This section turns to the families themselves: first whether the framework can count them, then whether it can tell them apart.

Read back on the step above, colour weakens at high energy for $n_{g} \leq 8$ and the weak force for $8n_{g} + n_{h} < 44$, so the running caps the number of generations without choosing it.

Neither consistency nor the running picks out three. The framework's own conjecture is that the number comes from counting: if the largest composite is three conformal spinors (twistors) and a massive unit uses two of them, there are \[ \binom{3}{2} = 3 \] ways to choose the pair. The bound of three is held only at observation level, so the conjecture is conditional on it.

\begin{conjecture}[Three generations as the faces of the maximal composite]\label{c:families.three-generations-faces-maximal-composite}
Matter comes in three copies, or generations, of each kind of particle, with no accepted explanation for why there are three. This conjecture proposes that the number comes from counting: if the largest allowed composite has three building blocks (a bound the series holds only at observation level) and a massive particle uses two of them, there are exactly three ways to choose the pair. Each choice would be one generation, with its lepton doublet as identified under Hypothesis (S), and the way generations mix would come from how the choices overlap.
\emph{First proved in} T13, Conj 3.1 (\href{https://doi.org/10.5281/zenodo.22722845}{doi:10.5281/zenodo.22722845}).
\end{conjecture}

The next step makes the counting precise. Each of the three faces carries an isomorphic copy of the fermion content, with the lepton doublet included under Hypothesis (S).

Any two faces share exactly one slot, which at observation level suggests a CKM-shaped mixing. That suggestion, and the matching of faces to real families, are not proved.

\begin{proposition}[Three identical families, mixing by shared slots]\label{c:families.three-identical-families-mixing-shared-slots}
If the largest allowed composite is made of three conformal spinors and a massive particle uses two of them, there are three ways to choose the pair. Each choice carries an exact copy of the same set of particles, because relabelling the spinors maps one copy onto another; the copies include the lepton doublet if the slots behave like the pairs (Hypothesis (S)). Any two choices share exactly one spinor, which hints at why quark families mix the way they do, though that hint and the matching of choices to real families are not proved.
\emph{First proved in} T13, Prop 3.2 (\href{https://doi.org/10.5281/zenodo.22722845}{doi:10.5281/zenodo.22722845}).
\end{proposition}
\noindent\emph{Rests on:} \ref{c:families.imported-one-derived-generation-hypothesis-s}, \ref{c:families.imported-face-valued-bracket-quantum-numbers}, \ref{c:standard-model.hyp-s-slot-statistics}.

Identical copies are not yet generations. In a quantum field theory three generations are three copies of the content under one gauge group, $3 \times 16 = 48$ independent states. The next step shows that the shared-slot picture of the step above cannot be read as three generations under one gauge group, so the mixing hint drawn from shared slots does not stand in that reading. If the three faces are read inside one state space, sharing their slots, the reading fails twice.

Two face planes in a three-dimensional slot space meet in a line, so any algebra acting on the slots leaves at most one face plane irreducible, and at most one face can be a weak doublet. The faces' chirality halves also overlap: they span $28$ of the $32$ even states with one pair system shared by the faces, and $46$, at most $47$, with separate pair modes per face.

\begin{theorem}[Read in one state space, the three faces of the three-spinor composite are not three generations under one gauge group: any Lie algebra acting linearly on the three slots has at most one face plane as an irreducible invariant subspace, so at most one face is a weak doublet, and the faces' chirality halves overlap: with one pair system shared by the faces they span 28 of the 32 even states of $\Lambda^{\bullet}(\C^{3} \oplus W)$ (at most 40 with halves chosen face by face), with separate pair modes per face 46 (at most 47), in every case short of the 48 of three generations]\label{c:synthesis.faces-of-one-slot-space-are-not-three-generations-under-one-gauge-group}
One proposal in the series says the three families of matter particles are the three ways of choosing two of the three building blocks of a largest composite. In a quantum field theory, three families means three separate copies of the particle content, all feeling the same forces in the same way. If the three choices live in one shared set of states, this fails twice: the weak force can treat only one of the three pairs as a doublet, because any two pairs share a block, and the shared blocks make the copies overlap, giving at most 28 distinct states (46 if each pair keeps its own colour modes) where three families need 48. The idea survives only if the three pairs are three separate composites.
\end{theorem}
\noindent\emph{Rests on:} \ref{c:families.three-identical-families-mixing-shared-slots}, \ref{c:standard-model.hyp-s-slot-statistics}, \ref{c:synthesis.rung-field-d-term-action-and-tree-masses-on-the-lift-locus}.

So the face identification survives only if the three faces are three separate composites, each with its own slots. The next step asks what the framework's standing premise on composites says about that reading.

The slot-module hypothesis, adopted as a premise, builds a composite's slot states from one vacuum by one doublet of slot operators, so every composite has exactly two slots. Under it the shared three-slot space fails even for the gauge group of a single face: either one slot goes uncounted by the centre that sets the hypercharge, or every species appears twice and the chirality half is vectorlike.

\begin{theorem}[Under the slot-module hypothesis the shared three-slot space of the faces is no slot module even for the gauge group of one face (a slot of weight zero, or two slot vacua with every species twice and a vectorlike chirality half), complementing the stored result that the faces are generations only as separate composites; the number of separate composites is independent of the premise and of anomaly freedom, and if both non-abelian factors are to be asymptotically free it is at most five at one or two doublets (both standard); that it is three is only an argued reading resting on the imported observation-level three-spinor bound, whose face space is the dual triplet of the slot space]\label{c:synthesis.slot-module-makes-the-faces-separate-composites-and-leaves-the-generation-number-free}
The framework's account of a particle family builds it from two internal ingredients called slots. Its standing assumption about how a family's states are organised already allows exactly two slots per particle, never three. This result shows that the old idea of the three families as the three pairs chosen from one three-slot object fails that assumption even when only one pair is looked at. In the shared space every particle would appear twice, and the left-right asymmetry of the weak force would cancel out. This adds to an earlier result that the families must be three separate objects. The same assumption, and the requirement that the theory be free of anomalies, are satisfied by any number of separate families, which is standard. If the strengths of both nuclear-type forces are to weaken at high energy, at most five families are allowed. That there are three is not derived: it is only a reading that rests on a separate, observation-level bound that the largest object has three slots, which can be paired in three ways.
\end{theorem}
\noindent\emph{Rests on:} \ref{c:synthesis.cpt-closure-forces-slot-module-and-chirality}, \ref{c:standard-model.imported-anomaly-coefficients}, \ref{c:synthesis.faces-of-one-slot-space-are-not-three-generations-under-one-gauge-group}, \ref{c:families.three-identical-families-mixing-shared-slots}, \ref{c:synthesis.one-loop-running-coefficients-from-spacetime-algebra-spin-traces}, \ref{c:standard-model.gauge-algebra-minimal-massive-composite}, \ref{c:synthesis.imported-witten-global-su2-anomaly}, \ref{c:standard-model.slot-module-hypothesis}, \ref{c:synthesis.krein-form-premise-indefinite-inner-product-for-gauge-fields}.

The step above is the plain answer to the counting question: the faces must be separate composites, and the series does not derive that there are three of them. The premise and anomaly cancellation hold copy by copy for any number of composites, and if both non-abelian forces are to weaken at high energy the running allows at most five at one or two Higgs doublets. Three appears only as a reading: the face space of three slots is three-dimensional, and taking it as the space of generation labels is a further family-symmetry hypothesis resting on the imported, observation-level three-spinor bound.

With the number taken as an input, the trail turns to what could make the identical families differ, starting with the vacuum.

Writing the deformed bracket as $B_{\Lambda} = b + \tfrac{\Lambda}{6}\,\Omega_{\omega}$, with $b$ the flat mass bracket, its square expands as \[ 2|B\_{\Lambda}|\^{}{2} = 2|b|\^{}{2} + \frac{2\Lambda}{3}\,$\mathrm{Re}$\bigl(\bar{b}\,\Omega\_{\omega}\bigr) + \frac{\Lambda\^{}{2}}{18}\,|\Omega\_{\omega}|\^{}{2}, \] whose middle term is linear in $b$. If the potential is built on this functional with $b$ and the vacuum $\omega$-area independent, a reading held at observation level, condensation cannot be avoided unless that area vanishes.

\begin{proposition}[The source term: the $\Lambda$-linear tilt of the mass functional]\label{c:families.source-term-linear-tilt-mass-functional}
Expanding the formula for a particle's mass with the cosmological constant switched on reveals an extra term that is linear in the basic mass quantity. Such a term acts like a slope that pushes the system away from zero. So, if the energy landscape is built from this formula and the two brackets in it are treated as independent, a nonzero value cannot be avoided unless the vacuum value of one of them, the $\omega$-area, vanishes. The size of that value then depends on that one quantity, which the framework names but does not compute. This second step is a reading at observation level, and the independence of the two brackets is the assumption it rests on.
\emph{First proved in} T13, Prop 4.2 (\href{https://doi.org/10.5281/zenodo.22722845}{doi:10.5281/zenodo.22722845}).
\end{proposition}
\noindent\emph{Rests on:} \ref{c:families.imported-dictionary-exact-join}, \ref{c:families.imported-face-valued-bracket-quantum-numbers}.

A natural guess was that this tilt also fixes the weak scale.

The tilt is $\Lambda$-suppressed, so for any vacuum length between the Planck and Hubble scales it lies between roughly $10^{-105}$ and $10^{-44}$ of the weak scale. The guess is refuted and kept on the record.

\begin{refutation}[The vacuum omega-area sets the weak scale]\label{c:refuted.the-vacuum-omega-area-sets-the-weak-scale}
It was a natural guess that the tiny vacuum area term which tilts the mass condensate also fixes the energy scale of the weak force. A size estimate shows it falls short by between forty-four and one hundred and five powers of ten, so the guess is wrong. What the vacuum term does fix is a direction among the three families, not a size, and the weak scale goes back to being an open question.
\end{refutation}

What the vacuum term can still fix is a direction among the families rather than a size, since directions are ratios and do not depend on $\Lambda$.

Under the indifference instrument that direction is democratic, and under the Gram-square reading it gives a rank-one mass form: one heavy family and two light ones at leading order. Separating the two light families is left as the next problem.

\begin{proposition}[The alignment computed: the thirteenth refutation, democracy, and one heavy family]\label{c:families.alignment-computed-thirteenth-refutation-democracy-one}
The paper first asks whether a tiny effect tied to the cosmological constant could set the energy scale of the weak force, and a size estimate shows it falls short by more than forty orders of magnitude, so that idea is ruled out. What the effect can still fix is a direction rather than a size. Assuming nature does not favour one family over another without cause, and assuming a particular form for the leading mass pattern, that direction gives one heavy family and two light ones at first approximation. Separating the two light families is left as the next problem.
\emph{First proved in} T13, Prop 4.3 (\href{https://doi.org/10.5281/zenodo.22722845}{doi:10.5281/zenodo.22722845}).
\end{proposition}
\noindent\emph{Rests on:} \ref{c:families.imported-indifference-instrument}, \ref{c:families.source-term-linear-tilt-mass-functional}.

With one direction singled out, the three identical faces can at last be ranked. Ordering them requires breaking $S_{3}$, the permutations of the three faces, which the condensate's alignment does at structural level.

As a first test the resulting nearest-neighbour texture gives the Cabibbo relation \[ \theta\_{C} \approx \sqrt{m\_{d}/m\_{s}} \approx 0.224, \] passed at observation level.

\begin{proposition}[The faces ordered: what the identification requires, and the first test passed]\label{c:families.faces-ordered-what-identification-requires-first}
The three candidate families in this framework start out exactly alike, but real families differ in mass, so something must tell them apart. The paper argues that the only thing available is the direction in which the vacuum settles, which then ranks the three by how much they overlap with it. As a first check, this kind of ranking leads to a well-known approximate formula linking a quark mixing angle to a ratio of quark masses, and the numbers agree to within about a percent. The check relies on standard assumptions borrowed from earlier work on such mass patterns.
\emph{First proved in} T13, Prop 4.4 (\href{https://doi.org/10.5281/zenodo.22722845}{doi:10.5281/zenodo.22722845}).
\end{proposition}
\noindent\emph{Rests on:} \ref{c:families.alignment-computed-thirteenth-refutation-democracy-one}, \ref{c:families.imported-face-valued-bracket-quantum-numbers}, \ref{c:families.three-identical-families-mixing-shared-slots}.

The next question is what the simplest vacua do to the two light families.

For every antisymmetric Gaussian pairing the mass form is \[ \cosh\^{}{2}r\,\delta\_{kl} + x(5 + 6x)\,\bar{n}\_{k}n\_{l}, \] rank one plus an ordering floor. So the light pair is exactly degenerate and $S_{2}$ cannot break in this class: whatever separates the electron and the muon must come from beyond it.

\begin{theorem}[The skeleton: rank one plus the ordering shift, and the unbreakable $S_{2}$]\label{c:families.skeleton-rank-one-plus-ordering-shift}
The paper computes exactly the mass pattern produced by its simplest vacuum, for every setting of its parameters. The answer is always one heavy family and two light ones with exactly equal masses, and the light masses come entirely from a quantum ordering effect. Because the two light masses can never differ in this setting, the electron and the muon cannot be told apart here, so whatever separates them must come from beyond this class of vacua.
\emph{First proved in} T13, Thm 6.2 (\href{https://doi.org/10.5281/zenodo.22722845}{doi:10.5281/zenodo.22722845}).
\end{theorem}
\noindent\emph{Rests on:} \ref{c:families.imported-squeezed-vacuum-bosonic-pairing}, \ref{c:families.rigidity-pairing-sector-one-event}.

The charged-lepton masses offer a sharp clue about what that missing ingredient must do. The linear conjecture predicts $Q_{m} = 2/3$, and the data give $0.8936$, so it is refuted.

In the $\sqrt{m}$ coordinate, suggested for the light pair by the rank-one leading form and assumed for the heavy family, \[ Q\_{\sqrt{m}} = \frac{\sum\_{k} m\_{k}}{\bigl(\sum\_{k}\sqrt{m\_{k}}\bigr)\^{}{2}} = 0.6666645, \] the Koide relation, holding within experimental uncertainty.

\begin{proposition}[The linear conjecture, refuted by data]\label{c:families.linear-conjecture-refuted-data}
A first guess, that the three charged-lepton masses themselves satisfy a simple balance condition, which they would if the masses were the basic coordinates and the split between the families were even, fails clearly against the data. The paper's own structure suggests why: the two lighter masses appear only at second order, so the natural quantities are their square roots. Using the square roots, the same condition holds to a few parts in a million. This is the Koide relation, known since the early 1980s; what the paper adds is a possible reason for using square roots, with an acknowledged gap for the heaviest lepton.
\emph{First proved in} T13, Prop 5.1 (\href{https://doi.org/10.5281/zenodo.22722845}{doi:10.5281/zenodo.22722845}).
\end{proposition}
\noindent\emph{Rests on:} \ref{c:families.alignment-computed-thirteenth-refutation-democracy-one}.

The obvious test is whether the simplest vacua can produce the Koide relation. Within the Gaussian class it holds only at the $S_{2}$ point, where it fixes a heavy-to-light ratio $K \approx 4616$.

The depth crosses the required ratio exactly once, at $x^{*} \approx K/6 \approx 769.33$, which is the Koide number transcribed. So the mechanism is a fit, and the derivation claim is refuted at this order.

\begin{proposition}[Koide in the class: the crossing, and the fit conclusion]\label{c:families.koide-class-crossing-fit-conclusion}
Within the paper's simplest class of vacuum states, the Koide mass relation can only hold in a special symmetric case, and there it fixes one number, the ratio of the heavy mass to the light ones. The model can match that number by choosing its depth parameter, and it does so at exactly one value. But that value is essentially the Koide number rewritten, so the model fits the relation rather than explaining it. The attempted derivation is therefore recorded as refuted at this level, and the relation is treated as a measurement of the depth.
\emph{First proved in} T13, Prop 6.3 (\href{https://doi.org/10.5281/zenodo.22722845}{doi:10.5281/zenodo.22722845}).
\end{proposition}
\noindent\emph{Rests on:} \ref{c:families.skeleton-rank-one-plus-ordering-shift}.

The failures add up to a single diagnosis.

At the displaced-Gaussian class the split, the angle and the depth are one open problem, since the simplest vacua take all three as inputs instead of predicting them. It is registered as a non-Gaussian completion that must produce the $207$ from caused structure.

\begin{conjecture}[One open problem in three forms]\label{c:families.one-open-problem-three-forms}
Three seemingly separate puzzles, the electron-muon mass ratio, the Koide angle and the depth parameter, turn out to be one unanswered question in this framework. Its simplest vacuum states take all three as inputs instead of predicting them. The paper registers the missing ingredient as its main open problem and sets conditions that any solution must satisfy.
\emph{First proved in} T13, Obs 7.1 (\href{https://doi.org/10.5281/zenodo.22722845}{doi:10.5281/zenodo.22722845}).
\end{conjecture}

Later work has sharpened that problem without solving it. Beyond the Gaussian class, on the states built from face words, the mass form is $H = (1 + \langle N\rangle/4)\mathbb{1}$ plus a positive weighted face density: a common floor plus a count of how the pairs are spread over the families.

\begin{theorem}[On the span of face words the mass form is the floor $1 + \langle N\rangle/4$ plus a weighted face density, so the forced eigenvector is the lowest in every source-built state and, under the mass-squared reading of the mass form, the split is a population ratio of the selected state]\label{c:synthesis.pair-singlet-mass-form-is-floor-plus-face-density}
In the families paper the three charged leptons get their masses from a vacuum full of correlated pairs of quanta, and the electron and muon come out equal unless the vacuum is of a more complicated, non-Gaussian kind. Earlier work showed that symmetry alone cannot fix the electron-to-muon ratio and noticed, in computed examples, that one particular direction always came out lightest. This result explains why, and does so for every vacuum of the relevant kind. The mass formula turns out to have a simple shape: a common floor, set only by how many pairs the vacuum holds on average, plus a term that counts how the pairs are spread over the family directions. The direction that always came out lightest is one that, by construction, never receives any pairs in these vacua, so it sits exactly on the floor. The other two masses then measure how the pairs are shared between the remaining two directions. Reading the formula's eigenvalues as squared masses, as the families paper does, the famous ratio of about 207 becomes a statement about how a frozen vacuum shares its pairs, and any rule that chooses the vacuum must get that sharing right.
\end{theorem}
\noindent\emph{Rests on:} \ref{c:families.imported-squeezed-vacuum-bosonic-pairing}, \ref{c:families.skeleton-rank-one-plus-ordering-shift}, \ref{c:synthesis.face-mass-form-unit-plus-gram}, \ref{c:synthesis.cumulants-at-fixed-second-moments-keep-the-floor}, \ref{c:synthesis.real-structure-breaks-source-stabiliser-to-z2}.

In the vacua built from the source and its conjugate one direction never receives pairs and sits exactly at the floor, which is why the lightest family always came out along it. Reading the eigenvalues as squared masses, as the families paper does, the electron's mass is fixed by the mean pair number and the ratio of about $207$ becomes a population ratio of the selected vacuum.

That reading sets a scale for any rule that picks the vacuum. For every state on the span of face words the heavier eigenvalues over the lightest are bounded by roughly twice the largest pair number, so the measured $(m_{\mu}^{2} + m_{\tau}^{2})/m_{e}^{2} \approx 1.213 \times 10^{7}$ needs weight at pair numbers of at least $6065574$.

\begin{theorem}[Under the mass-squared reading of the face mass form, every state on the span of face words has $(\lambda_{2} + \lambda_{3})/\lambda_{1} \le (2 + \tfrac52\langle D\rangle + \langle D^{2}\rangle)/(1 + \tfrac12\langle D\rangle)$, and a state supported on pair numbers at most $K$ has it at most $2K + 1 + 2/(K + 2)$, so the measured charged-lepton masses at the pole masses need a frozen vacuum with weight at pair numbers of at least $6065574$ (an estimated $6.4$ million at common-scale running masses, first order in $\alpha$)]\label{c:synthesis.mass-squared-reading-needs-millions-of-face-pairs}
In one reading of the theory, the three charged-lepton masses squared are the three levels of a matrix built from the vacuum, and the vacuum is described by how many pairs of certain quanta it holds. This result shows that the ratio of the heavy levels to the light one can never exceed about twice the largest number of pairs present. Since the tau is about 3477 times heavier than the electron (at the pole masses), the squared ratio is about twelve million, so the vacuum must contain states with at least about six million pairs; with the masses defined at a common energy scale instead, the number rises to about 6.4 million. Whatever rule picks the vacuum has to produce a number that large from the theory's structure, without it being put in by hand.
\end{theorem}
\noindent\emph{Rests on:} \ref{c:synthesis.pair-singlet-mass-form-is-floor-plus-face-density}, \ref{c:synthesis.koide-channel-phase-target-depends-on-mass-definition}, \ref{c:synthesis.imported-one-loop-qed-pole-mass-relation}.

\section{Results used from elsewhere}
The proofs above use the following results of the series, stated briefly; each is proved in the paper named.

\begin{fact}[Imported: the dictionary and the exact join]\label{c:families.imported-dictionary-exact-join}
Mass is weight read against the frozen point, $m^{2} = -w(w+3)H^{2}$, and for two conformal spinors $m^{2} = 2|B_{\Lambda}|^{2} + \Lambda C/3$, which puts a heavy composite's mass in the imaginary part of its weight.
\emph{Source:} T13, Fact 2.2 (\href{https://doi.org/10.5281/zenodo.22722845}{doi:10.5281/zenodo.22722845}).
\end{fact}

\begin{fact}[Imported: the face-valued bracket and its quantum numbers]\label{c:families.imported-face-valued-bracket-quantum-numbers}
The mass bracket transforms by $\det(U)$ under the internal $U(2)$ and carries the quantum numbers $(1,1)$ with $Y = +1$ of $e^{c}$, not those of a weak doublet.
\emph{Source:} T13, Fact 2.1 (\href{https://doi.org/10.5281/zenodo.22722845}{doi:10.5281/zenodo.22722845}).
\end{fact}

\begin{fact}[Imported: the indifference instrument]\label{c:families.imported-indifference-instrument}
On the families paper's reading, indifference, a structural premise, equates weights only along symmetry orbits and never equates inequivalent irreducible components.
\emph{Source:} T13, Fact 2.5 (\href{https://doi.org/10.5281/zenodo.22722845}{doi:10.5281/zenodo.22722845}).
\end{fact}

\begin{fact}[Imported: one derived generation, and Hypothesis (S)]\label{c:families.imported-one-derived-generation-hypothesis-s}
The quark doublet and the singlet hypercharges follow unconditionally, while the lepton doublet, and, given the one-handed content and $Y_{Q} = 1/6$, the anomaly-fixed uniqueness of the generation's hypercharges, hold only under Hypothesis (S).
\emph{Source:} T13, Fact 2.3 (\href{https://doi.org/10.5281/zenodo.22722845}{doi:10.5281/zenodo.22722845}).
\end{fact}

\begin{fact}[Imported: the squeezed vacuum of a bosonic pairing]\label{c:families.imported-squeezed-vacuum-bosonic-pairing}
The squeezed vacuum of a bosonic pairing $\xi = U\Sigma U^{\mathsf{T}}$ has two-point functions $N = \overline{U\sinh^{2}\Sigma\,U^{\dagger}}$ and $M = -U\sinh\Sigma\cosh\Sigma\,U^{\mathsf{T}}$, and Wick's theorem gives its four-point functions from them.
\emph{Source:} T13, Fact 2.6 (\href{https://doi.org/10.5281/zenodo.22722845}{doi:10.5281/zenodo.22722845}).
\end{fact}

\begin{fact}[Rigidity: the pairing sector is one event]\label{c:families.rigidity-pairing-sector-one-event}
Every antisymmetric pairing has family factor $[a]_{\times}$ with $GG^{\dagger}$ spectrum $\{|a|^{2}, |a|^{2}, 0\}$, so the freezing is a single fourfold-degenerate Bogoliubov event of depth $r = \tau|a|$.
\emph{Source:} T13, Prop 6.1 (\href{https://doi.org/10.5281/zenodo.22722845}{doi:10.5281/zenodo.22722845}).
\end{fact}

\begin{fact}[The Fock space and the Maxwell modes]\label{c:field-theory.fock-space-maxwell-modes}
With $i = I_4$ the photon's Fock space is the polynomial CCR representation with one $I_4$-complex mode per null momentum, and positive frequency is positive helicity for every massless even field.
\emph{Source:} FT, Prop 8.1.
\end{fact}

\begin{fact}[Second quantisation]\label{c:field-theory.second-quantisation}
With the single-algebra field and $i = I_4$, second quantisation gives a local field and the standard $d\Gamma(M)$, one chirality entering complex-linearly and the other conjugated.
\emph{Source:} FT, Prop 5.2.
\end{fact}

\begin{fact}[The gauge algebra of the minimal massive composite]\label{c:standard-model.gauge-algebra-minimal-massive-composite}
The automorphism algebra of the minimal massive hierarchy $\C^{2} \otimes \C^{3}$ is $\su(2) \oplus \su(3) \oplus \uu(1)$, exactly the Standard Model's, while that of $\C^{3} \otimes \C^{3}$ is seventeen-dimensional.
\emph{Source:} T12, Thm 4.1 (\href{https://doi.org/10.5281/zenodo.22722669}{doi:10.5281/zenodo.22722669}).
\end{fact}

\begin{fact}[Hypothesis (S): slot statistics, derived from the slot-module hypothesis]\label{c:standard-model.hyp-s-slot-statistics}
Under the slot-module hypothesis, the composite's slots are fermionic modes obeying exclusion and forming one Clifford system with its pairs, its state space is the five-mode Fock space $\Lambda^{\bullet}(\C^{2} \oplus W)$ with $U(2) \times U(3)$ acting as $\Lambda^{\bullet}(A \oplus B)$, chirality is the parity of the mode number, and slots of one tower parity whose grading is the chirality are odd.
\emph{Source:} T12, Hypothesis (S) (\href{https://doi.org/10.5281/zenodo.22722669}{doi:10.5281/zenodo.22722669}).
\end{fact}

\begin{fact}[Hypercharge from the global form, pair counting, and consistency]\label{c:standard-model.hypercharge-global-form-pair-counting-consistency}
Under Hypothesis (S) for the doublets, the composite's states carry the observed hypercharges of one generation, which satisfy all four anomaly conditions and are the unique solution for that content with $Y_{Q} = \tfrac16$.
\emph{Source:} T12, Thm 5.5 (\href{https://doi.org/10.5281/zenodo.22722669}{doi:10.5281/zenodo.22722669}).
\end{fact}

\begin{fact}[Imported: the anomaly coefficients]\label{c:standard-model.imported-anomaly-coefficients}
A $U(1)$ gauge factor acting on Weyl fermions is consistent only if four anomaly sums vanish: the two mixed with $SU(3)$ and $SU(2)$, the cubic one and the gravitational one.
\emph{Source:} S. Weinberg, The Quantum Theory of Fields, Vol. 2: Modern Applications, Cambridge University Press (1996)., \S 22.4.
\end{fact}

\begin{fact}[Imported: one-loop running, and the trace ratio of a simple unifying group]\label{c:standard-model.imported-one-loop-running-trace-ratio}
At one loop the inverse couplings run linearly in $\ln\mu$ with the stated coefficients, and if a simple unifying group exists they meet in the normalisation fixed by its trace ratios.
\emph{Source:} H. Georgi and S. L. Glashow, Unity of all elementary-particle forces, Phys. Rev. Lett. 32 (1974), 438--441.; U. Amaldi, W. de Boer, and H. F"urstenau, Comparison of grand unified theories with electroweak and strong coupling constants measured at LEP, Phys. Lett. B 260 (1991), 447--455.; S. Navas et al. (Particle Data Group), Review of particle physics, Phys. Rev. D 110 (2024), 030001..
\end{fact}

\begin{fact}[Anomaly freedom forces odd slot parity and the slot content of Hypothesis (S), up to multiplicity]\label{c:synthesis.anomaly-freedom-forces-slot-content}
With the composite space a $U(2)$ slot module tensored with the pair system, hypercharge from the centre and chirality the tower grading $(-1)^{\pi n + k}$, the particle half is anomaly-free if and only if the slots are odd ($\pi = 1$) and the slot module is $m$ copies of $\Lambda^{\bullet}\C^{2}$, and then equivariant slot creation obeys exclusion; parafermionic slots of order $p \ge 2$ are anomalous.
\emph{Source:} Transport Lab claim synthesis.anomaly-freedom-forces-slot-content, version 3.
\end{fact}

\begin{fact}[For a slot module generated from one vacuum by one doublet of slot creation operators, CPT closure forces the fermionic slot module and anomaly freedom forces the chirality to be the mode-number parity away from the neutral singlets]\label{c:synthesis.cpt-closure-forces-slot-module-and-chirality}
For $V = F \otimes \Lambda^{\bullet}W$ with $Y = (3n - 2k)/6$, CPT closure holds iff $F = a\Lambda^{0} \oplus b\Lambda^{1} \oplus a\Lambda^{2}$; with one vacuum anomaly-free halvings need $b$ odd, and with $F$ generated from its vacuum CPT closure forces $F = \Lambda^{\bullet}\C^{2}$ with exclusion and anomaly freedom forces the chirality $\pm(-1)^{n + k}$ off the neutral singlets.
\emph{Source:} Transport Lab claim synthesis.cpt-closure-forces-slot-module-and-chirality, version 1.
\end{fact}

\begin{fact}[Higher cumulants added at the skeleton's second moments cannot lower the light floor]\label{c:synthesis.cumulants-at-fixed-second-moments-keep-the-floor}
The number part of the face mass form is $\tfrac12(n\mathbb{1} - \rho^{\mathsf{T}})$ (one-body family density), the skeleton floor $\cosh^{2}r$ saturates the bound it gives, and any state with the skeleton's second moments has $H \ge \cosh^{2}r\,\mathbb{1}$: cumulants only raise masses, so reaching $m_{\mu}/m_{e} = 207$ at the Koide depth needs the connected Gram matrix's top eigenvalue about $7{,}900$ times the Gaussian one and a new second eigenvalue above $3 \times 10^{7}$.
\emph{Source:} Transport Lab claim synthesis.cumulants-at-fixed-second-moments-keep-the-floor, version 1.
\end{fact}

\begin{fact}[In every state the face mass form is at least the unit, the excess an explicit Gram matrix]\label{c:synthesis.face-mass-form-unit-plus-gram}
For every state of the six bosonic face modes the mass form is $H = \mathbb{1} + W^{\dagger}W$ with $W$ a state-dependent linear map from face space, so $H \ge \mathbb{1}$ in every state and every family splitting lies in the excess $W^{\dagger}W$.
\emph{Source:} Transport Lab claim synthesis.face-mass-form-unit-plus-gram, version 1.
\end{fact}

\begin{fact}[Imported: the axial anomaly of spinor electrodynamics (Adler 1969; Bell and Jackiw 1969): when the vector current is kept conserved, the divergence of the axial current of a Dirac fermion has, besides the mass term, the anomalous term $(\alpha/4\pi)\,\epsilon^{\xi\sigma\tau\rho}F_{\xi\sigma}F_{\tau\rho}$ from the triangle graph, which accounts for the decay of the neutral pion into two photons]\label{c:synthesis.imported-adler-bell-jackiw-axial-anomaly}
The axial anomaly (Adler; Bell and Jackiw 1969): with the vector current conserved, $\partial\cdot j^{5} = 2im\,j^{5} + (\alpha/4\pi)\epsilon FF$ from the triangle graph, proportional to $E\cdot B$; it explains $\pi^{0} \to \gamma\gamma$.
\emph{Source:} S. L. Adler, Axial-vector vertex in spinor electrodynamics, Phys. Rev. 177 (1969), 2426--2438.; J. S. Bell and R. Jackiw, A PCAC puzzle: pi0 -> gamma gamma in the sigma-model, Nuovo Cimento A 60 (1969), 47--61.; S. L. Adler, Anomalies to all orders, in: 50 Years of Yang-Mills Theory, ed. G. 't Hooft, World Scientific (2005); arXiv:hep-th/0405040., eq. (1) and the reference list.
\end{fact}

\begin{fact}[Imported: the heat-kernel (Seeley-DeWitt) coefficients of a Laplace-type operator, whose field-strength part on a flat background is $\tfrac{1}{12}\Omega_{\mu\nu}\Omega^{\mu\nu} + \tfrac12E^{2}$ ($a_{2}$ in DeWitt's numbering, $a_{4}$ in Vassilevich's), and the background-field gauge with Faddeev-Popov ghosts and the one-loop determinant weights (Abbott 1981)]\label{c:synthesis.imported-heat-kernel-coefficients-and-background-field-gauge}
The heat-kernel expansion $\mathrm{Tr}\,e^{-tD} \simeq \sum t^{(k-n)/2}a_{k}$, with flat-background field-strength coefficient $\tfrac12E^{2} + \tfrac1{12}\Omega^{2}$ ($a_{4}$ in Vassilevich's labelling, $a_{2}$ in DeWitt's) multiplying the logarithm in $\ln\det$, and the background-field gauge with Faddeev-Popov ghosts and the one-loop determinant weights (Abbott 1981).
\emph{Source:} D. V. Vassilevich, Heat kernel expansion: user's manual, Phys. Rep. 388 (2003), 279--360, arXiv:hep-th/0306138.; B. S. DeWitt, Dynamical Theory of Groups and Fields, Gordon and Breach (1965).; L. F. Abbott, The background field method beyond one loop, Nucl. Phys. B 185 (1981), 189--203..
\end{fact}

\begin{fact}[Imported: the one-loop machinery of covariant perturbation theory: Feynman parametrisation, the Wick-rotated integrals of $(\ell^{2} - \Delta)^{-n}$ with symmetric integration, Pauli-Villars regulator fields of either sign of weight (Pauli and Villars 1949), the covariant photon propagator with gauge-parameter independence of on-shell amplitudes, and on-shell external states]\label{c:synthesis.imported-one-loop-integrals-and-pauli-villars-regulators}
Feynman parametrisation, the Wick-rotated integrals ($-i/(32\pi^{2}\Delta)$ for three denominators, $\tfrac{i}{16\pi^{2}}(\ln(L/\Delta) - 1)$ for two) with symmetric integration, Pauli-Villars regulators with $\sum c_{i} = \sum c_{i}M_{i}^{2} = 0$ and weights of both signs, the covariant photon propagator with gauge-parameter independence of on-shell amplitudes, and on-shell external states.
\emph{Source:} W. Pauli and F. Villars, On the invariant regularization in relativistic quantum theory, Rev. Mod. Phys. 21 (1949), 434--444.; M. E. Peskin and D. V. Schroeder, An Introduction to Quantum Field Theory, Addison-Wesley (1995), chapters 6 and 7.; S. Weinzierl, Introduction to Feynman integrals, arXiv:1005.1855, equations (50) and (61)..
\end{fact}

\begin{fact}[Imported: at one loop in QED the pole mass of a charged lepton is m\_pole = [1 + (alpha/pi)(1 + (3/4) ln(mu\^{}2/mbar\^{}2))] mbar(mu), and Koide's relation written in running masses is shifted by about 0.1 percent (Sumino 2009)]\label{c:synthesis.imported-one-loop-qed-pole-mass-relation}
$m^{\mathrm{pole}} = [1 + (\alpha/\pi)(1 + \tfrac34\ln(\mu^{2}/\bar{m}^{2}))]\bar{m}(\mu)$ at one loop in QED, so running-mass ratios differ from pole-mass ratios at order $\alpha$ (Sumino 2009).
\emph{Source:} Y. Sumino, Family gauge symmetry as an origin of Koide's mass formula and charged lepton spectrum, arXiv:0903.3640 (2009), report TU-842., eq. (6); Y. Sumino, Family gauge symmetry as an origin of Koide's mass formula and charged lepton spectrum, J. High Energy Phys. 05 (2009), 075..
\end{fact}

\begin{fact}[Imported: Witten's global $SU(2)$ anomaly; a four-dimensional $SU(2)$ gauge theory with an odd number of left-handed fermion doublets (and no other representations) is inconsistent, so a consistent theory has an even number of them (Witten 1982)]\label{c:synthesis.imported-witten-global-su2-anomaly}
An $SU(2)$ gauge theory with an odd number of left-handed fermion doublets is inconsistent ($\pi_{4}(SU(2)) = \mathbb{Z}_{2}$ flips the sign of the fermion integral), so the number of doublets must be even (Witten 1982).
\emph{Source:} E. Witten, An SU(2) anomaly, Phys. Lett. B 117 (1982), no. 5, 324--328..
\end{fact}

\begin{fact}[In the $\sqrt{m}$ (circulant) reading the channel-phase target is a pole-mass statement: at first order in $\alpha$ the one-loop QED relation between pole and $\overline{\mathrm{MS}}$ masses moves the Koide phase reproducing $m_{\mu}/m_{e}$ at a common scale to $(0.22254, 0.22258)$, more than $1800$ times the distance of $\tfrac29$ from $\delta_{*}$, and puts Koide's $Q$ off $\tfrac23$ by $(1.0, 1.3) \times 10^{-3}$]\label{c:synthesis.koide-channel-phase-target-depends-on-mass-definition}
In the $\sqrt{m}$ (circulant) reading, at first order in $\alpha$, the one-loop QED relation between pole and $\overline{\mathrm{MS}}$ masses moves the Koide phase that reproduces $m_{\mu}/m_{e}$ at a common scale from $\delta_{*} \approx 0.2222220$ to $(0.22254, 0.22258)$ (more than $1800$ times the distance of $\tfrac29$) and puts Koide's $Q$ off $\tfrac23$ by $(1.0, 1.3) \times 10^{-3}$, so the stored nine-decimal target and its symmetry-angle bound $4208$ (which falls to $61$) are pole-mass statements; under the mass-squared reading only the ratio shift and the Koide defect carry over.
\emph{Source:} Transport Lab claim synthesis.koide-channel-phase-target-depends-on-mass-definition, version 2.
\end{fact}

\begin{fact}[One-loop electrodynamics in the spacetime algebra, in a covariant gauge with imported Pauli-Villars regulators: transverse vacuum polarisation with coefficient $e^{2}/(12\pi^{2})$, $F_{2}(0) = \alpha/(2\pi)$ through the real bivector $\gamma^{\mu} \wedge q$, $Z_{1} = Z_{2}$, and the unit outside the product on massive lines]\label{c:synthesis.one-loop-qed-in-the-spacetime-algebra}
Under the Krein-form premise, in a covariant gauge with imported Pauli-Villars regulators (some of negative weight), one-loop electrodynamics in the spacetime algebra gives a transverse vacuum polarisation with coefficient $e^{2}/(12\pi^{2})$, $F_{2}(0) = \alpha/(2\pi)$ through the real bivector $\gamma^{\mu} \wedge q$, and $Z_{1} = Z_{2}$, with the unit outside the product on massive lines.
\emph{Source:} Transport Lab claim synthesis.one-loop-qed-in-the-spacetime-algebra, version 2.
\end{fact}

\begin{fact}[Every $S_{3}$-equivariant real structure breaks the coherent source's stabiliser to $\mathbb{Z}_{2}$, forcing exactly one eigenvector of the mass form (the reflection of the Koide vector through the democratic axis; at the floor, so the lowest, in every state built from the source, its conjugate and the invariant quartic word, and not in general) and leaving a five-parameter block that the state alone fixes; the channel phase is invisible to it and a cubic family word forces nothing, so the $207$ cannot come from the real structure or the cubic family word acting by covariance]\label{c:synthesis.real-structure-breaks-source-stabiliser-to-z2}
Every $S_{3}$-equivariant real structure and the coherent source $\kappa(u' + \lambda e^{i\chi}f_{+})$ are carried by a face unitary to conjugation and the line of $(1, a, ia)$, $a = \lambda/\sqrt2$, with a phase independent of $\chi$; the residual symmetry is the single reflection $R_{n}$, $n = c_{R}\times c_{I}$ (at the Koide point the reflection of the Koide vector through the democratic axis), which forces $n$ to be an exact eigenvector of the face mass form (at the floor, so the lowest, in every state built from the source, its conjugate and $\hat{V}$, not in general) and leaves a five-parameter block that exact Fock instances fill, while a cubic family word has no residual symmetry at all: the $207$ cannot come from the real structure or the cubic family word acting by covariance, only from the selected state.
\emph{Source:} Transport Lab claim synthesis.real-structure-breaks-source-stabiliser-to-z2, version 2.
\end{fact}

\begin{fact}[On the null-lift locus every $U(2)$-invariant potential of the rung field is a function of the Hermitian norm $N$ and the hypercharge moment map $m$ on the orbit space $N^{2} - m^{2} \ge H^{4}$, whose boundary is the charge-preserving lift family; the hypercharge D-term action puts the vacuum on that boundary at the level $m^{*} = g'^{2}\xi/(g^{2} + g'^{2})$, so $v/H$ enters only through the input $\xi$, and gives the neutral mode the mass $M_{Z}$ and the charged pair $M_{W}$ at tree level in every frame; an added invariant $W(N)$ shifts them to $M_{Z}^{2} + 2W''m^{*2}/N^{*} + 2W'H^{4}/N^{*2}$ and $M_{W}^{2} + 2W'$, the latter twice the outward derivative of the potential, so local stability against charge breaking is the positivity of the charged mass; the rung field carries a physical charged scalar pair besides the neutral mode]\label{c:synthesis.rung-field-d-term-action-and-tree-masses-on-the-lift-locus}
On the null-lift locus the rung field's invariant potentials are functions of $(N, m)$ on $N^{2} - m^{2} \ge H^{4}$ with the lift family as boundary; the hypercharge D-term action has a charge-preserving vacuum at the level $m^{*} = g'^{2}\xi/(g^{2} + g'^{2})$, so $v/H$ is set by the input $\xi$ (the premise registering the level as the free number; under the promotion $\xi G$, and so $v/M_{P}$, is the law and $v/H$ initial data), with $m_{h} = M_{Z}$ and $m_{\pm} = M_{W}$ at tree level in every frame; an added $W(N)$ gives $m_{h}^{2} = M_{Z}^{2} + 2W''m^{*2}/N^{*} + 2W'H^{4}/N^{*2}$ and $m_{\pm}^{2} = M_{W}^{2} + 2W'$, twice the outward derivative, so local stability against charge breaking is the positivity of the charged mass, and the field carries a physical charged pair.
\emph{Source:} Transport Lab claim synthesis.rung-field-d-term-action-and-tree-masses-on-the-lift-locus, version 3.
\end{fact}

\section*{Conclusion}
On the count, the trail ends with a negative answer: the series does not derive three generations. Anomaly cancellation holds generation by generation for any number, the running allows at most five if both non-abelian forces are to weaken at high energy, and three rests on an imported observation-level bound read through a family-symmetry hypothesis. The matching of the three faces to the real families is itself still a conjecture, now read as three separate composites.

On the masses, three identical families become one heavy family and two light ones once the vacuum picks a direction, while separating the light pair remains open. Under the mass-squared reading any vacuum that does it must hold millions of face pairs, a number some rule has to produce from the theory's structure rather than by hand.

Several steps rest on stated premises: the slot-module hypothesis and Hypothesis (S) that it yields, the Krein inner product for gauge fields in the running, the observation-level three-spinor bound, and the indifference premise that nature does not favour one family without cause.

\end{document}